\section*{Abstract}

%Magnetic vortices have gotten significant research attention in recent years and are increasingly regarded as promising candidates for next-generation non-volatile memory technologies. In this work, we investigate how the geometry and size of polygonal magnetic micro- and nanostructures affect their \textcolor{red}{magnetic domains, curl density, and whether or not a vortex is formed}. Using STEM-DPC measurements of the lateral magnetic field distribution, we calculated the z-component of the curl for a selection of polygonal geometries and normalized this value by the polygon area to enable direct comparison across different sizes. \textcolor{red}{The results show that polygons with fewer sides tend to have higher magnetic anisotropy, while polygons with seven or more sides increasingly stabilize into vortices with uniform curl-density. Moreover, the normalized curl drops off sharply between five and six sides, but gradually flattens out for polygons with higher symmetry. These observations suggest that the emergence of vortex states is primarily driven by shape-induced magnetic anisotropy, which diminishes as the geometry approaches circular symmetry.}

Magnetic domain behavior in nanostructures is influenced by geometry and size. In this report, we investigate how different polygonal micro- and nanostructures affects the formation of magnetic domains in permalloy nanostructures fabricated with FIB. Using STEM-DPC, we analyzed the curl of the in-plane magnetization for polygons ranging from triangles to nonagons with radii of approximately 400, 600, and 900 nm. The results show a clear transition from discrete magnetic domains in low-sided polygons to more continuous configurations as the number of sides increases. The average curl follows an exponential decay with increasing symmetry, consistent across all size scales. These findings demonstrate that the transition from discrete to continuous magnetic domains is primarily geometry-driven, highlighting the role of polygonal shape in stabilizing vortex states in nanomagnetic systems.
